\documentclass[11pt]{article}
\usepackage[margin=1in]{geometry}
\usepackage{amsmath,amssymb,amsthm,mathtools}
\usepackage[hidelinks]{hyperref}
\usepackage{microtype}

\newtheorem{theorem}{Theorem}
\newtheorem{lemma}{Lemma}
\newtheorem{remark}{Remark}
\newcommand{\BMO}{\mathrm{BMO}}
\newcommand{\R}{\mathbb R}
\newcommand{\1}{\mathbf 1}

\title{Distinct-symbol bilinear commutators:\\
an affirmative answer for the sum and a counterexample for the iterate}
\author{Run \texttt{fa\_banach\_001}, lane 7}
\date{11 August 2026}

\begin{document}
\maketitle

\begin{abstract}
We resolve the two final questions of arXiv:1707.01639 in the standard
real-symbol setting of the paper.  For a homogeneous, nondegenerate bilinear
convolution Calder\'on--Zygmund operator, boundedness of
\([b_1,T]_1+[b_2,T]_2\) forces \(b_1,b_2\in\BMO\).  The proof first detects
\(b_1+b_2\), then subtracts its ordinary commutator and detects \(b_1\) from
the difference of the two slots.  In contrast, the iterated-commutator
question is false as printed: a constant inner symbol annihilates the entire
iterate, while the other symbol may be outside BMO.
\end{abstract}

\section{Setting and result}

Let
\[
 T(f,g)(x)=\iint_{\R^n\times\R^n}K(x-y_1,x-y_2)f(y_1)g(y_2)\,dy_1dy_2
\]
off the common support.  We use the kernel hypotheses intended in the source:
\(K\) is a continuous bilinear Calder\'on--Zygmund kernel, homogeneous of
degree \(-2n\), and nonzero away from the origin in \(\R^{2n}\).  The last
condition follows in the source from the assumption that \(1/K\) has an
absolutely convergent Fourier series on every cube avoiding the origin.  We
also use the standard sufficiency theorem, proved in the source, that
\(a\in\BMO\) implies boundedness of each \([a,T]_j\) in the full bilinear
range.

For locally integrable symbols set
\[
 \mathcal S_{b_1,b_2}=[b_1,T]_1+[b_2,T]_2,
 \qquad
 \mathcal P_{b_1,b_2}=[b_2,[b_1,T]_1]_2.
\]

\begin{theorem}[Split resolution of Problems A and B]
\label{thm:main}
Let \(b_1,b_2\) be real-valued and locally integrable, let
\(1<p_1,p_2<\infty\), and put
\(1/q=1/p_1+1/p_2\) (so \(q>1/2\)).
\begin{enumerate}
\item If
\(\mathcal S_{b_1,b_2}:L^{p_1}\times L^{p_2}\to L^q\) is bounded, then
\(b_1,b_2\in\BMO\).
\item Boundedness of \(\mathcal P_{b_1,b_2}\) does not imply
\((b_1,b_2)\in\BMO\times\BMO\), even if the two symbols are distinct.
Indeed, one may take \(b_1=0\) and \(b_2(x)=x_1\).
\end{enumerate}
The first assertion also holds for complex symbols whenever the kernel is
real-valued up to a fixed scalar phase.  For arbitrary complex kernels, the
same proof is complete for real symbols; the source's questions and examples
are normally read in this real-symbol setting.
\end{theorem}

\section{A separated-cube sector lemma}

\begin{lemma}[Uniform sector]
\label{lem:sector}
There are a vector \(z\in\R^n\), a fixed number \(R>1\), a unimodular
\(\sigma\), and \(c_K>0\) with the following property.  If \(P\) is a cube
of side length \(\ell\), and \(Q=P+R\ell z\), then
\[
 \operatorname{Re}\!\left(\sigma K(x-y_1,x-y_2)\right)
       \ge c_K\ell^{-2n}
 \quad(x\in Q,\ y_1,y_2\in P).
\]
The roles of a prescribed cube and its translate may be reversed.
\end{lemma}

\begin{proof}
Choose \(z\ne0\).  Nonvanishing gives \(K(z,z)\ne0\); rotate it by a
unimodular \(\sigma\) so that \(\sigma K(z,z)>0\).  If \(x\in Q\) and
\(y_i\in P\), then
\[
 \ell^{2n}K(x-y_1,x-y_2)
 =K(Rz+u_1,Rz+u_2)
 =R^{-2n}K(z+u_1/R,z+u_2/R),
\]
where \(|u_i|\le C_n\).  Continuity at \((z,z)\) makes the last expression
stay in a fixed open half-plane when \(R\) is sufficiently large.  Since
\(R\) is then fixed, the asserted lower bound follows.  Translating in the
opposite direction gives the reversed geometry.
\end{proof}

\section{Problem A: recovery of both symbols}

Write \(s=b_1+b_2\), and let \(A\) denote the operator norm of
\(\mathcal S_{b_1,b_2}\).

\subsection*{Step 1: recover the sum}

Fix an arbitrary cube \(Q\) of side \(\ell\), and use Lemma
\ref{lem:sector} to choose a separated cube \(P\) of the same size on which
the two inputs will be supported.  Let \(m_i\) be a median of \(b_i\) on
\(P\), and define
\[
 E_i^- =\{y\in P:b_i(y)\le m_i\},\qquad
 E_i^+ =\{y\in P:b_i(y)\ge m_i\}.
\]
Both sets have measure at least \(|P|/2\).  Put \(m=m_1+m_2\).  The kernel
formula for the summed commutator is
\[
 \mathcal S_{b_1,b_2}(f,g)(x)
 =\iint K(x-y_1,x-y_2)
 \bigl(s(x)-b_1(y_1)-b_2(y_2)\bigr)f(y_1)g(y_2)\,dy_1dy_2.
\]
For \(x\in Q\cap\{s\ge m\}\) and \(y_i\in E_i^-\), the parenthesis is a
nonnegative real number at least \(s(x)-m\).  Lemma \ref{lem:sector} therefore
gives
\[
 (s(x)-m)_+
 \le \frac{C\ell^{2n}}{|E_1^-||E_2^-|}
 \left|\mathcal S_{b_1,b_2}(\1_{E_1^-},\1_{E_2^-})(x)\right|.
\]
Taking the \(L^q(Q)\) quasi-norm and using boundedness yields
\[
 \|(s-m)_+\|_{L^q(Q)}
 \le C A\ell^{2n}|E_1^-|^{1/p_1-1}|E_2^-|^{1/p_2-1}
 \le CA|Q|^{1/q}.
\]
The upper median sets give the same estimate for \((s-m)_-\).  Hence
\[
 \sup_Q\inf_{c\in\R}
 \left(\frac1{|Q|}\int_Q|s-c|^q\right)^{1/q}\le CA.
\]
The \(L^q\)-oscillation characterization of BMO is valid for every
\(q>0\) (for \(q<1\), this is precisely the \(\BMO_q=\BMO\) lemma proved in
the target paper).  Thus \(s\in\BMO\).

\subsection*{Step 2: recover one symbol from the slot difference}

Since \(s\in\BMO\), the ordinary commutator \([s,T]_2\) is bounded in the
same range.  Consequently
\[
 \mathcal D_{b_1}:=\mathcal S_{b_1,b_2}-[s,T]_2
 =[b_1,T]_1-[b_1,T]_2
\]
is bounded.  Its kernel identity has the crucial cancellation
\[
 \mathcal D_{b_1}(f,g)(x)
 =\iint K(x-y_1,x-y_2)
       \bigl(b_1(y_2)-b_1(y_1)\bigr)f(y_1)g(y_2)\,dy_1dy_2.       \tag{1}
\]

Now fix an arbitrary cube \(P\), let \(m\) be a median of \(b_1\) there,
and set
\[
 E=\{b_1\le m\}\cap P,\qquad F=\{b_1\ge m\}\cap P.
\]
Again \(|E|,|F|\ge|P|/2\).  Choose the separated output cube \(Q\) from
Lemma \ref{lem:sector}.  On \(E\times F\), the factor in (1) is nonnegative,
so for every \(x\in Q\),
\[
 |\mathcal D_{b_1}(\1_E,\1_F)(x)|
 \ge c_K\ell^{-2n} I_P,
 \qquad
 I_P:=\int_E\int_F(b_1(y_2)-b_1(y_1))\,dy_2dy_1.
\]
The operator bound, \(|Q|=|P|\), and
\(1/q=1/p_1+1/p_2\) imply
\[
 c_K\ell^{-2n}I_P|Q|^{1/q}
 \le C|E|^{1/p_1}|F|^{1/p_2}
 \le C|P|^{1/q},
\]
and hence \(I_P\le C|P|^2\).  On the other hand,
\[
 I_P
 =|E|\int_F(b_1-m)+|F|\int_E(m-b_1)
 \ge \frac{|P|}{2}\int_P|b_1-m|.
\]
Therefore
\[
 \frac1{|P|}\int_P|b_1-m|\le C
\]
uniformly in \(P\), proving \(b_1\in\BMO\).  Finally
\(b_2=s-b_1\in\BMO\).  This proves Theorem \ref{thm:main}(1).

\section{Problem B: a universal counterexample}

Take
\[
 b_1(x)=0,\qquad b_2(x)=x_1.
\]
The symbols are distinct, \(b_1\in\BMO\), and \(b_2\notin\BMO\), since its
mean oscillation on cubes of side \(L\) grows proportionally to \(L\).
Nevertheless
\[
 [b_1,T]_1=0,
 \qquad
 \mathcal P_{b_1,b_2}=[b_2,[b_1,T]_1]_2=[b_2,0]_2=0.
\]
The zero bilinear operator is bounded between every pair of quasi-Banach
function spaces.  Thus Problem B has a negative answer under every kernel and
exponent hypothesis in the source.  The obstruction is intrinsic: an
iterated commutator only sees the joint oscillation of the tuple and cannot
force each factor into BMO without an individual nondegeneracy assumption.

\section{Scope and checks}

The affirmative argument uses exactly two structural inputs beyond boundedness
of \(T\): a fixed-phase separated-cube lower bound, supplied by homogeneity and
nonvanishing, and the already-known BMO sufficiency for ordinary commutators.
It does not use the Fourier-series representation of \(1/K\).

For complex-valued symbols and a real kernel (or a kernel with one fixed
global phase), apply the real argument to real and imaginary parts, using
complex conjugation and linearity.  A variable-phase complex kernel with
complex symbols requires a finite-sector version of the median lemma; that
corner is not claimed here.  Problem B's counterexample has no such
restriction.

\begin{thebibliography}{9}
\bibitem{target}
D. Wang, J. Zhou, and Z. Teng,
\emph{Characterizations of weighted BMO space and its application},
arXiv:1707.01639.
\end{thebibliography}

\end{document}
